\section{Details on the theoretical dipole}\label{app:b-1M-derivation}
Here we compute the kinematic and clustering dipoles
in terms of the harmonic transfer functions,
$\Delta\Zu{\kin}_\ell$ and $\Delta\Zu{\clu}_\ell$, respectively.
First, we recall that for a general field $O(\n)$ the coefficients
$a_{\ell m}$ can always be expressed as a $\bk$-space integral over the primordial perturbations,
weighted by the appropriate harmonic transfer function $\Delta_\ell(k)$:
\be
    a\Z{\ell m}
        =(-\im)^\ell\sp4\pi\int\frac{\dif^3\bk}{(2\pi)^3}\:
            Y^*_{\ell m}(\hat\bk)\sp \Delta_\ell(k)\sp {\zeta}(\bk).
    \label{eqn:a-1M}
    \ee
Then given the transfer functions for any two observables, labelled $A$ and $B$,
the spectrum $C_\ell^{AB}$ is given by
\be
    C_\ell^{AB}
        = 4\pi\int^\infty_0\frac{\dif k}{k}\:
        	\Delta\Zu{A}_\ell(k)\sp \Delta\Zu{B}_\ell(k)\sp\PR(k),
    \label{eqn:C-ell-ab}
\ee
where we assume a standard primordial (adiabatic) power spectrum $\PR(k)$ defined by
$\langle{\zeta}(\bk){\zeta}^*(\bk')\rangle=(2\pi)^3\delta_{\rm D}(\bk-\bk')\sp 2\pi^2\PR(k)/k^3$. In principle, superhorizon isocurvature fluctuations can also be considered, which are known to give rise to an
intrinsic dipole in the CMB~\citep{Turner:1991dn}. However, it was
recently shown~\citep{Domenech:2022mvt} that such modes cannot also
induce a sizeable intrinsic dipole in the number counts.

Now, using equations~\eqref{eqn:a-1M} and \eqref{eqn:dipole-vecs}, we find that the model
prediction for the dipoles $\mbf{d}_\mrm{X}$, $\mrm{X}\in\{\kin,\clu\}$, can be written as
\be
    \mbf{d}_\mrm{X}
        = -\ssp4\pi\im\sqrt{\frac{3}{4\pi}}
          \int\frac{\dif^3\bk}{(2\pi)^3}\, \mat{U}^\dagger\ssp\bm{Y}_1(\hat\bk)\ssp\Delta\Zu{\mrm{X}}_1(k)\ssp {\zeta}(\bk),
\ee
where $\bm{Y}_1\equiv(Y_{1-1},Y_{11},Y_{10})\Zu\T$ and
we have the unitary matrix
\be\label{eq:U}
    \mathbf{U}
    =
    \begin{pmatrix}
        \!\phantom{-}1/\sqrt{\ssp2}    &   \im/\sqrt{\ssp2}  &   0   \\[2pt]
        \!          -1/\sqrt{\ssp2}    &   \im/\sqrt{\ssp2}  &   0   \\[2pt]
        \!0           &       0        &   1
    \end{pmatrix},
\ee
with $\mat{U}\mat{U}^\dagger=\mat{U}^\dagger\mat{U}=\mat{I}_3$.
Note that $\mat{U}^\dagger\ssp\bm{Y}_1$ is a vector containing
the $\ell=1$ real-valued spherical harmonics.

Computing the statistics of
$\mbf{d}_\mrm{X}$ -- i.e.~$\langle{d}_{\mrm{X}i}{d}_{\mrm{Y}j}\rangle=\sigma^2\Z{\mrm{XY}}\sp \delta_{ij}^\mrm{K}$ -- is now straightforward.
To obtain the dispersion $\sigma^2\Z{\mrm{XY}}$ we can
simply select the $z$-component, all components being statistically equivalent
under isotropy. Using orthogonality of spherical harmonics,
we have
%% $\sigma\Z{\mrm{XY}}^2
%% =\langle d\Z{\mrm{X}}^{\sp x}\sp d\Z{\mrm{Y}}^{\sp x}\rangle
%% =\langle d\Z{\mrm{X}}^{\sp y}\sp d\Z{\mrm{Y}}^{\sp y}\rangle
%% =\langle d\Z{\mrm{X}}^{\sp z}\sp d\Z{\mrm{Y}}^{\sp z}\rangle$
\bea
    \sigma\Z{\mrm{XY}}^2
        % &=(4\pi)^2\ssp\frac{3}{4\pi}\int\frac{\dif^2\hat\bk}{4\pi}\
        %     Y_{10}(\hat\bk)\ssp Y_{10}\Zu{*}(\hat\bk) \nonumber\\
        % &\quad\times\int^\infty_0\frac{k^2\dif k}{2\pi^2}\,
        % \Delta\Zu{\mrm{X}}_1(k)\ssp\Delta\Zu{\mrm{Y}}_1(k) \ssp \frac{2\pi^2}{k^3}\ssp\PR(k) \nonumber\\[2pt]
        &=3\int^\infty_0\frac{\dif k}{k}\ \Delta_1\Zu{\mrm{X}}(k)\ssp \Delta_1\Zu{\mrm{Y}}(k)\ssp \PR(k)
        =\frac{3}{4\pi}\ssp C\Zu{\mrm{XY}}\Z1,
    %% \qquad
    %% \mrm{X},\mrm{Y} \in \{\kin,\clu\}.
    \label{eqn:sig-ab}
\eea
where we have used
$\langle{\zeta}(\bk)\ssp{\zeta}^*(\bk')\rangle=(2\pi)^3\delta_{\rm D}(\bk-\bk')\sp 2\pi^2\PR(k)/k^3$.

Note that in this appendix the $a_{\ell m}$'s belong to the full sky.
Those belonging to the cut sky can always be obtained
in terms of the full-sky $a_{\ell m}$'s by geometric 
linear relations.
The explicit expressions are given later
in Appendix~\ref{app:mode-coupling}.


\subsection{Clustering transfer function} 
Clustering gives rise to fluctuations $\delta_\qso(\n)$ in the number counts. In spherical harmonics,
\be
	\delta_\qso(\n)
	=\sum_{\ell,m}^\infty a\Z{\ell m}\, Y_{\ell m}(\n),
\quad
	a\Z{\ell m}=\int\dif^2\n~Y_{\ell m}^*(\n)\ssp\delta_\qso(\n).
\ee
Here we are interested in the dipole $\dint\cdot\n=\sum_{m=-1}^1 a_{1m}\sp Y_{1m}(\n)$.
Substituting $\delta_\qso(\n)=\int\dif z\,p_z(z)\ssp\delta_\qso(\n,z)$, with $p_z$ the
normalised redshift distribution, and the Fourier
representation of $\delta_\qso(\n,z)$, we have
\begin{align*}
a\Z{\ell m}
&=\int\dif^2\n~Y_{\ell m}^*(\n)\int^\infty_0\dif z~p_z(z)\ssp\delta_\qso(\n,z) \nonumber\\[-2pt]
&=\int\dif^2\n~Y_{\ell m}^*(\n)\int^\infty_0\dif z~p_z(z)\int\frac{\dif^3\bk}{(2\pi)^3}~
	{\delta_\qso}(\bk,z)\ssp \rme\Zu{-\im\bk\cdot\chi(z)\n},
\end{align*}
where in the second line we write $\x=\chi(z)\n$.
Now substituting the plane-wave expansion,
$\rme^{-\im\bk\cdot\chi\n}
=\sum_{\ell,m}(-\im)^{\ell}4\pi\sp j_{\ell}(k\chi)\sp
Y^*_{\ell m}(\hat\bk)\sp Y^{\phantom*}\Z{\ell m}(\n)$,
%\be\label{eqn:plane-wave-exp}
%	\rme\Zu{-\im\bk\cdot\chi(z)\n}
%        =\sum_{\ell,m}(-\im)^{\ell}4\pi\sp j_{\ell}\big(k\chi(z)\big)\sp
%         Y^*_{\ell m}(\hat\bk)\sp Y^{\phantom*}\Z{\ell m}(\n)
%\ee
into the foregoing expression we find
\begin{align*}
a\Z{\ell m}
&=\int\dif^2\n~Y_{\ell m}^*(\n)
	\int^\infty_0\dif z~p_z(z)\int\frac{\dif^3\bk}{(2\pi)^3}~{\delta_\qso}(\bk,z) \nonumber\\
	&\qquad\qquad\qquad\times\sum_{\ell',m'}(-\im)^{\ell'}4\pi\ssp j_{\ell'}\big(k\chi(z)\big)\ssp
	Y^*_{\ell'm'}(\hat{\bk})\ssp Y_{\ell'm'}(\n) \nonumber\\
&=(-\im)^\ell\sp 4\pi\int\frac{\dif^3\bk}{(2\pi)^3}\ Y^*_{\ell m}(\hat{\bk}) \nonumber\\
	&\qquad\:\times\bigg[\int^\infty_0\dif z~p_z(z)\ssp \bqso(z)\ssp\mathcal{T}_\delta(k,z)\ssp
			j_\ell\big(k\chi(z)\big)+\cdots\bigg]\zeta(\bk),
\end{align*}
where in the second line we used the
orthogonality of spherical harmonics.
Here we have inserted
$\delta_\qso(\bk,z)=\bqso(z)\ssp\delta(\bk,z)+\cdots=\bqso(z)\ssp\mathcal{T}_\delta(k,z)\ssp\zeta(\bk)+\cdots$,
where the ellipsis are corrections from 
redshift-space distortions, Doppler effects, and general-relativistic
corrections 
\citep[see, e.g.,][appendix A, for complete expressions]{CLASSgal}.
We thus have
\be
	\Delta\Zu{\clu}_\ell(k)
	=\int^\infty_0\dif z~p_z(z)\ssp\bqso(z)\ssp\mathcal{T}_\delta(k,z)\ssp
			j_\ell\big(k\chi(z)\big)
   +\cdots.
\ee
% In this work we use the \citet{2005MNRAS.356..415C}
% quasar bias parametrisation, $\bqso(z)=0.53+0.289(1+z)^2$.


\subsection{Kinematic transfer function}\label{app:calc-b1m}
For the kinematic dipole recall that $\dkin\cdot\n=\ckin\,\bm\beta\cdot\n$, with
$\ckin=2+x(1+\alpha)$ and $\bm\beta=\v_O/c=\v(\x\!=\!0)/c$. In harmonics,
\be
	\dkin\cdot\n
	=\sum_{L,M} b\Z{LM}\, Y_{LM}(\n),
\quad
	b\Z{LM}=\int\dif^2\n~Y_{LM}^*(\n)\ssp\dkin\cdot\n,
\ee
the calculation carries through in much the same way as with the previous calculation for number
counts, {provided} that we keep $\x$ arbitrary, setting $\x=0$ only at the end of the
calculation. (Note that since we are dealing with a pure dipole only
$L=1$ coefficients are nonzero.) We thus write equation~\eqref{eqn:b-LM} as
\bea
    b_{1M}
        &=\int\dif^2\n\ Y_{1M}^*(\n)\ssp\ckin\,\v\Z{R}(\x\!=\!0)\cdot\n/c \nonumber\\[-2pt]
        &=\int\dif^2\n\ Y_{1M}^*(\n)\ssp\ckin
            \int\frac{\dif^3\bk}{(2\pi)^3}
            \frac{-\im\bk\cdot\n}{k}\bigg(\frac{H_0 f_0}{c\sp k}\bigg)\sp W(kR)\sp\delta(\bk)
    \label{eqn:b-1M-w1}
\eea
and then use that with $\x=\chi\n$, where $\chi$ the comoving distance, the dipole 
$-\im\sp\bk\cdot\n$ can be expressed as 
$-\im\sp\bk\cdot\n=(\partial\sp\rme^{-\im\bk\cdot\chi\n}/\partial\chi)|_{\chi=0}$.
Now substituting the plane-wave expansion
%\be
%    -\im\k\cdot\n/k=\frac1k\frac{\partial}{\partial\chi}\rme^{-\im\bk\cdot\chi\n}
%        =\frac1k\frac{\partial}{\partial\chi}
%         \bigg[\sum_{\ell,m}(-\im)^{\ell}4\pi\sp j_{\ell}(k\chi)\sp
%         Y^*_{\ell m}(\hat\bk)\sp Y^{\phantom*}\Z{\ell m}(\n)\bigg]_{\chi=0},
%\ee
into equation~\eqref{eqn:b-1M-w1} we find
\begin{align*}
    b_{1M}
        &=\int\dif^2\n\ Y_{1M}^*(\n)\sp\ckin
            \int\frac{\dif^3\bk}{(2\pi)^3}
            \bigg(\frac{H_0f_0}{c\sp k}\bigg) W(kR)\sp\delta(\bk) \nonumber\\
        &\qquad\qquad\times
            \frac1k\frac{\partial}{\partial\chi}
            \bigg[\sum_{\ell,m}(-\im)^{\ell}4\pi\sp j_\ell(k\chi)\sp
                    Y^*_{\ell m}(\hat\bk)\sp Y\Z{\ell m}(\n)\bigg]_{\chi=0} \nonumber\\[2pt]
        &=(-\im)\sp4\pi\int\frac{\dif^3\bk}{(2\pi)^3}\:
            Y^*_{1M}(\hat\bk)\sp \ckin\bigg(\frac{H_0f_0}{c\sp k}\bigg)W(kR)\sp\delta(\bk)
                \frac1k\frac{\partial j_1}{\partial\chi}\bigg|_{\chi=0},
\end{align*}
where in the second line we have used the orthogonality of spherical
harmonics.
Since $j_1(x)=x/3+\mathcal{O}(x^3)$, the derivative with respect to $\chi$, at $\chi=0$,
exactly evaluates to $k/3$; inserting $\delta(\bk)=\mathcal{T}_\delta(k,z)\sp{\zeta}(\bk)$, we finally get
\be\nonumber
    b_{1M}
        =(-\im)\sp4\pi\int\frac{\dif^3\bk}{(2\pi)^3}\:
            Y_{1M}^*(\hat\bk)
                \bigg[\frac13\sp\ckin\sp W(kR)\bigg(\frac{H_0f_0}{c\sp k}\bigg)\sp \mathcal{T}_\delta(k)\bigg]\sp{\zeta}(\bk),
\ee
where $\mathcal{T}_\delta(k)$ is evaluated at $z=0$. Note that
$\bqso$ does not appear here because $\v$ is sourced by
matter perturbations. Since this expression is in the form of equation~\eqref{eqn:a-1M} we can immediately
read off the kinematic transfer function from
the contents of the square brackets:
\be
    \Delta\Zu{\kin}_1(k)
        =\frac13\sp\ckin\sp W(kR)\sp\bigg(\frac{H_0f_0}{c\sp k}\bigg)\sp\mathcal{T}_\delta(k).
        %% =-\ckin\sp W(kR)\sp\frac{\mathcal{T}_\theta(k)}{3\sp c\sp k}
    \label{eqn:Delta-kin-1}
\ee
Note that because we evaluate $\v$ at $\x=0$ the 
transfer function $\mathcal{T}_\delta(k)$ is simply evaluated
at $z=0$ and not integrated. Further note that to
include source evolution,
simply replace $\ckin$ with its average $\bar{A}_\kin$ 
given by equation~\eqref{eqn:d-kin-corr}.







\section{Pixelization}
We can also consider the effect of the pixelization on the theoretical power.
At the field level, the pixelization of the number-count fluctuations is given by
\be
    \delta\Z{\qso,i} = \int\dif^2\n\: W^\mrm{pix}_i(\n)\ssp \dqso(\n),
\ee
where $W^\mrm{pix}_i(\n)$ is the \healpix\ pixel window function for the $i$th pixel,
which is equal to zero, unless $\n$ falls within the $i$th pixel, in which case
$W^\mrm{pix}_i$ is equal to $1/\Omega_{\rm pix}$, with $\Omega_{\rm pix}$ being the
pixel area. This window function is normalised,$\int\dif^2\n\,W^\mrm{pix}_i(\n)=1$,
so that $\delta\Z{\qso,i}$ is the average fluctuation in the $i$th pixel.
Note that because the \healpix\ pixel shape varies azimuthally there is no global pixel window function
that applies to all pixels. Provided a large enough $N_\mrm{side}$ is chosen this is not really a problem: the differences between pixel window functions
only becomes important when considering large $\ell$, but this can always be remedied by choosing
a larger $N_\mrm{side}$.

We will ignore the azimuthal variation in pixels,
as is routinely done. Under this
approximation pixelized fields are also statistically isotropic and the angular power spectrum
of the pixelized field is
$
    C^\mrm{\ssp pix}_\ell = (\sp\bar{W}^\mrm{\ssp pix}_\ell\sp)^2 C_\ell
$,
where $C_\ell$ is the unpixelized power, and $\bar{W}^\mrm{\sp pix}_\ell$ is the \healpix
pixel-averaged window function~\citep{2005ApJ...622..759G}. Since $\bar{W}^\mrm{\sp pix}_\ell\leq1$ there is a loss of power for all
$\ell$ (more for larger $\ell$).
In this work $N_\mrm{side}=64$ with which we find that
$\bar{W}^\mrm{\ssp pix}_\ell$ is equal to unity to
within $0.5\%$ for $\ell=0$ to $\ell=20$. As this power loss is negligible we have thus
ignored the effects of pixelization on the theoretical power.
% taking $C^\mrm{\ssp pix}_\ell=C_\ell$.


% First thirty W_ell for nside=64
%% [1.        , 0.99997725, 0.99993174, 0.99986349, 0.9997725 ,
%%  0.99965876, 0.99952229, 0.9993631 , 0.99918118, 0.99897657,
%%  0.99874926, 0.99849926, 0.9982266 , 0.99793129, 0.99761334,
%%  0.99727278, 0.99690962, 0.99652388, 0.99611559, 0.99568476,
%%  0.99523142, 0.99475561, 0.99425734, 0.99373664, 0.99319354,
%%  0.99262808, 0.99204028, 0.99143018, 0.99079781, 0.99014322]



\section{Derivation of $\Pr(\Dtot)$}\label{app:pD_derivation}
In this appendix we derive the probability distribution function $p(\Dtot)$ in both the full-sky and cut-sky regimes. The PDF in the former case is well known; it is the Maxwellian. In the latter case, considering an arbitrary symmetric Galactic
plane cut (the details of which will be fixed in Appendix~\ref{app:gal-cut}),
we derive another long-tailed PDF, one that is suppressed relative to the Maxwellian.
The calculation comes down to the marginalisation~\eqref{eqn:p-Dtot-int}, which we evaluate through the moment generating function. The full-sky PDF is parametrised by $\sigkin$,
$\sigint$, and $\sigkinint$; whilst the cut-sky PDF is parametrised by
$\sigma_{\kin_x\kin_x}$, $\sigma_{\kin_z\kin_z}$,
$\sigma_{\clu_x\clu_x}$, $\sigma_{\clu_z\clu_z}$, $\sigma_{\kin_x\clu_x}$, and
$\sigma_{\kin_z\clu_z}$ (these are, however, linearly related to those of the full sky).
The calculation is based on geometrical considerations;
no cosmological model needs to be assumed.

\subsection{Full sky}\label{app:Dtot-details}
To set up the problem, construct the six-dimensional vector $\bm{X}=(\dkin,\dint)\Zu\T$,
with $\bm{X}\sim\calN(\mbf{0},\cov)$.
By statistical isotropy only components along the same axis can
be correlated ($x$-$x$, $y$-$y$, etc).
Thus $\cov$ can be written in block matrix form~\eqref{eqn:cov-6x6}, with each block 
proportional to the $3\times3$ identity matrix,
$\langle\mathbf{d}_\mrm{X}\Zu{\phantom\T} \mathbf{d}_\mrm{Y}\Zu\T\rangle\propto\mathbf{I}_3$,
with $\mrm{X},\mrm{Y}\in\{k,c\}$.
In particular, we have that the expected amplitude of the kinematic dipole is
$\Dkin\equiv\langle\dkin^2\rangle^{1/2}=\sqrt{\sp3}\ssp\sigkin$,
and the expected amplitude of the clustering dipole is
$\Dint\equiv\langle\dint^2\rangle^{1/2}=\sqrt{\sp3}\ssp\sigint$,
where the root three is due to $\sigkin$ and $\sigint$ being
one-dimensional dispersions, while $\Dkin$ and $\Dint$ is the total dispersion. Evaluating $\sigkin$, $\sigint$,
and $\sigkinint$ requires specifying a cosmological model, but these are simply related
to the $\ell=1$ multipole of certain angular power spectra by $\sigma\Z{\mrm{XY}}\Zu2=3C^\mrm{XY}_1/(4\pi)$.


Though $\cov$ is a $6\times6$ matrix there are only three parameters needed to specify it (on
account of isotropy). In analysing this matrix we find it convenient to therefore
write it as a Kronecker product:
\be
    \cov
        =
        \begin{pmatrix}
            \,\sigkin^2 & \sigkinint^2\, \\[3pt]
            \,\sigkinint^2 & \sigint^2\,
        \end{pmatrix}
        \otimes \mathbf{I}_3,
    \qquad
    \cov^{-1}
        =
        \begin{pmatrix}
            \,\sigkin^2 & \sigkinint^2\, \\[3pt]
            \,\sigkinint^2 & \sigint^2\,
        \end{pmatrix}^{-1}
        \otimes \mathbf{I}_3,
\ee
where $\mathbf{I}_3$ is the $3\times3$ identity matrix, and the inverse $\cov^{-1}$ is
the Kronecker product of the inverse of the two matrices.%
\footnote{$\mat{A}^{-1}=(\mat{B}\otimes\mat{C})^{-1}=\mat{B}^{-1}\otimes\mat{C}^{-1}$.}
Thus we see that the covariances in $\cov$ are essentially contained in a $2\times2$ matrix.


The square of the amplitude of the total dipole, $Y\equiv\Dtot^2=|\dkin+\dint|^2$, can be
recast as the quadratic form
\be
    Y=\bm{X}\Zu\T\mathbf{A}\bm{X},
    \qquad
    \mathbf{A}
    \equiv
        \begin{pmatrix}
            \,\mathbf{I}_3 & \mathbf{I}_3\, \\
            \,\mathbf{I}_3 & \mathbf{I}_3\,
        \end{pmatrix}
        =
        \begin{pmatrix}
            \sp1 & 1\, \\
            \sp1 & 1\,
        \end{pmatrix}
        \otimes\mathbf{I}_3,
        % \begin{pmatrix}
        %     1 & 0 & 0 & 2 & 0 & 0 \\[2pt]
        %     0 & 1 & 0 & 0 & 2 & 0 \\[2pt]
        %     0 & 0 & 1 & 0 & 0 & 2 \\[2pt]
        %     2 & 0 & 0 & 1 & 0 & 0 \\[2pt]
        %     0 & 2 & 0 & 0 & 1 & 0 \\[2pt]
        %     0 & 0 & 2 & 0 & 0 & 1
        % \end{pmatrix},
\ee
where $\mathbf{A}$ is a $6\times6$ symmetric matrix.
The moment generating function
${M}\Z{Y}(t)\equiv\langle \rme^{tY}\rangle$ can be written using the
`law of the unconscious statistician' as
\be
    {M}\Z{Y}(t)
        \equiv\int\dif Y\, \rme^{tY} p(Y)
        =\int\dif^6\bm{X}\, \rme^{t\bm{X}\Zu\T\mathbf{A}\bm{X}} p(\bm{X}),
    \label{eqn:M-Y}
\ee
i.e.~in terms of the probability distribution of $\bm{X}$.
Before we compute the second integral, we recall that ${M}_Y(-t)$ is the
(unilateral) Laplace transform
$\mathscr{L}$ of the PDF $p(Y)$:
\be
        \mathscr{L}[p(Y)](t)
        \equiv\int^\infty_0 \dif Y\, \rme^{-tY} p(Y)
        ={M}\Z{Y}(-t).
\ee
Therefore, once the moment generating function is in hand, and the sign of its argument
is flipped, we recover the PDF by simply taking the inverse Laplace transform:
\be
    p(Y)
        =\mathscr{L}^{-1}\big[{M}\Z{Y}(-t)\big](Y)
        \equiv\frac{1}{2\pi\im}\int^{\gamma+\im\infty}_{\gamma-\im\infty}\dif t\,
            \rme^{tY}{M}\Z{Y}(-t).
\ee
(Here $\gamma$ is some large, positive constant, chosen
so that the poles lie to the
left of it in the complex plane.)
In practice, we will not have to perform such integrals as we can build the solution
out of simpler, known transforms, by using the convolution properties of the Laplace transform.%
\footnote{Perhaps the more obvious approach to obtain the PDF is to use the characteristic function
$\phi\Z{Y}(t)\equiv{M}\Z{Y}(\im t)$ together with the Fourier transform.
However, this proves slightly awkward, since $Y\geq0$ forces us to consider a
one-sided Fourier transform.}


To evaluate the second integral in equation~\eqref{eqn:M-Y}, we note that it has the basic form
of a Gaussian convolution; because one `Gaussian' is without its normalisation factor, 
we find after some straightforward matrix algebra
\be
    {M}\Z{Y}(t)
        =\det\big(\mathbf{I}_6 - 2\sp t\mathbf{A}\cov\big)^{-1/2}
        =\prod_{j=1}^6 \lambdaup\Z{j}^{-1/2},
\ee
where in the second equality we have expressed the determinant in terms of the
eigenvalues $\lambdaup_j$ of $\mathbf{I}_6-2\sp t\mathbf{A}\cov$. Now specialising to
the case of statistical isotropy, we find that there are only two
distinct eigenvalues, which we call $\lambdaup_1$ and $\lambdaup_2$.
To see this we exploit the properties of the Kronecker product: write
$\mathbf{I}_6=\mathbf{I}_2\otimes\mathbf{I}_3$ and
$\mathbf{A}\cov=\bm\Sigma\otimes\mathbf{I}_3$, with
\be
    \bm\Sigma
        \equiv
        %% \begingroup % keep the change local
        %% \setlength\arraycolsep{2.5pt}
        \begin{pmatrix}
            \sp1 & 1\, \\
            \sp1 & 1\,
        \end{pmatrix}
        \begin{pmatrix}
            \sp\sigkin^2 & \sigkinint^2\, \\[3pt]
            \sp\sigkinint^2 & \sigint^2\,
        \end{pmatrix}
        %% \endgroup
        =
        \begin{pmatrix}
            \,\sigkin^2+\sigkinint^2 & \sigint^2+\sigkinint^2\, \\[4pt]
            \,\sigkin^2+\sigkinint^2 & \sigint^2+\sigkinint^2\,
        \end{pmatrix}.
\ee
Because
$\mathbf{I}_6-2\sp t\mathbf{A}\cov
    =(\mathbf{I}_2-2\sp t\bm\Sigma)\otimes\mathbf{I}_3$, there are two distinct eigenvalues
(each having multiplicity three) corresponding to the $2\times2$ matrix $\mathbf{I}_2-2\sp t\bm\Sigma$.
The eigenvalues are $\lambdaup_1=1$ and $\lambdaup_2=1-2\sp t\Delta^2$,
with
\be
    \Delta^2
        \equiv\sigkin^2 + 2\sp\rho\sp\sigkin\sp\sigint + \sigint^2.
\ee
% Note that $\Delta^2$ is bounded below when rho=-1 and so $\Delta^2\geq(\sigkin-\sigint)^2$.
Therefore, we have simply that ${M}_Y(t)=(1-2\sp t\Delta^2)^{-3/2}$, which
can be seen to closely resemble the moment generating function of the generalised
chi-squared distribution.
As mentioned, the PDF is recovered by inverse Laplace transform of ${M}_Y(-t)$:%
\footnote{The integral can be done with help of the convolution theorem, which
says that a product of moment generating functions corresponds to the convolution of their
individual PDFs.
Thus we write $(1-2\sp t\Delta^2\big)^{-3/2}=\prod_{j=1}^3(1-2\sp t\Delta^2\big)^{-1/2}$,
and note that the chi-squared distribution with one degree of freedom has a moment
generating function $(1-2\sp t\Delta^2\big)^{-1/2}$. Then the PDF
corresponding to ${M}_Y(t)$ is the convolution $p(Y)=p_x\ast p_y\ast p_z$ of
three such chi-squared distributions, which yields another chi-squared distribution
but with three degrees of freedom.}
\be
    p(Y)
        =\frac{1}{2\pi\im}\int^{\gamma+\im\infty}_{\gamma-\im\infty}
            \frac{\rme^{tY}\dif t}{\big(1+2\sp t\Delta^2\big)^{3/2}}
        = \frac{\sqrt{\sp Y}\sp \rme^{-Y/(2\Delta^2)}}{\sqrt{\sp2\pi}\ssp\Delta^3}.
\ee
Finally, the distribution for the amplitude $\Dtot\equiv\sqrt{\sp Y}$ is
${p(\Dtot)=p(Y\!=\!\Dtot^2)\ssp|\sp\dif Y/\dif\Dtot|}$, or, explicitly, after some algebra,
\be
    p(\Dtot)\ssp\dif\Dtot
        =\frac{4}{\sqrt\pi}\ssp\exp\bigg(-\frac{\Dtot^2}{2\Delta^2}\bigg)\ssp
        \frac{\Dtot\Zu2}{2\sp\Delta\Zu2}\sp\frac{\dif \Dtot}{\sqrt2\sp\sp\Delta}.
    \label{eqn:p-X}
\ee
This PDF, a chi distribution with three degrees of freedom (or Maxwellian), 
has mean $\langle\Dtot\rangle=\sqrt{\sp8/\pi\sp}\sp\Delta$ and
variance $\langle\Dtot^2\rangle=3\Delta^2$. However, the variance
is larger than might be expected from a standard Maxwellian due to cross terms
($\sigma^2_{\kin_x\clu_x}$, $\sigma^2_{\kin_z\clu_z}$, etc).

\begin{figure}
\centering
    \begin{subfigure}{\columnwidth}
        \centering
	    \includegraphics[width=\textwidth]{figures/Figure7} % p_fullsky_with_source_evoln.pdf
    \end{subfigure}
    \caption{As for Fig.~\ref{fig:p-Dtot-cutsky} but on the full sky ($f_\mrm{sky}=1$).
    The probability distribution is simply the Maxwellian~\eqref{eqn:p-X}
    with dispersions
    given by equations~\eqref{eq:fullsky_dispersions}. The main point here is
    that without the mask larger values of $\Dtot$ are more probable compared
    with the case of including the mask}.
    \label{fig:p-Dtot}
\end{figure}


\subsection{Cut sky}\label{app:p-Dtot-cutsky}
The Galactic plane cut, described later in Appendix~\ref{app:galactic-cut}, has the effect of identifying a preferred direction, i.e.~along the $z$-axis (where the $z$-axis is orthogonal to the Galactic plane).
Because of this the covariance now
takes the form
\bea\label{eqn:cov-6x6-cutsky}
    \cov
    % &=
    %     \left(
    %     \begin{array}{cc}
    %         \big\langle\dkin\Zu{\phantom\T}\dkin\Zu\T\big\rangle &
    %         \big\langle\dkin\Zu{\phantom\T}\dint\Zu\T\big\rangle \\[7pt]
    %         \big\langle\dkin\Zu{\phantom\T}\dint\Zu\T\big\rangle &
    %         \big\langle\dint\Zu{\phantom\T}\dint\Zu\T\big\rangle
    %     \end{array}\right) \nonumber\\[5pt]
    &=
        \left(
        \begin{array}{cccccc}
            \sigma^2_{\kin_x\kin_x} & 0 & 0 & \,\sigma^2_{\kin_x\clu_x} & 0 & 0 \\
            0 & \sigma^2_{\kin_x\kin_x} & 0  & \,0 & \sigma^2_{\kin_x\clu_x} & 0 \\
            0 & 0 & \sigma^2_{\kin_z\kin_z} & \,0 & 0 & \sigma^2_{\kin_z\clu_z} \\
            \sigma^2_{\kin_x\clu_x} & 0 & 0\ & \sigma^2_{\clu_x\clu_x} & 0 & 0 \\
            0 & \sigma^2_{\kin_x\clu_x} & 0\ & 0 & \sigma^2_{\clu_x\clu_x} & 0 \\
            0 & 0 & \sigma^2_{\kin_z\clu_z} & 0 & 0 & \sigma^2_{\clu_z\clu_z}
        \end{array}\right).
\eea
Notice now that each $3\times3$ block matrix is no longer proportional to the identity matrix,
so cannot be written as a Kronecker product.
In particular, we now see that the dispersions of the $x$- and
$y$-components differ from that of the $z$-component. Because of this
$\sigma^2_{\clu_x\clu_x}=\sigma^2_{\clu_y\clu_y}\neq\sigma^2_{\clu_z\clu_z}$ and
$\sigma^2_{\kin_x\kin_x}=\sigma^2_{\kin_y\kin_y}\neq\sigma^2_{\kin_z\kin_z}$; similarly,
for the cross-terms
$\sigma^2_{\kin_x\clu_x}=\sigma^2_{\kin_y\clu_y}\neq\sigma^2_{\kin_z\clu_z}$.
The probability distribution $p(\Dtot)$ is now given in terms of six independent quantities,
$\sigma^2_{\kin_x\kin_x}$, $\sigma^2_{\kin_z\kin_z}$,
$\sigma^2_{\clu_x\clu_x}$, $\sigma^2_{\clu_z\clu_z}$, $\sigma^2_{\kin_x\clu_x}$, and $\sigma^2_{\kin_z\clu_z}$.


We still have azimuthal
symmetry (rotations about the $z$-axis), and this is reflected in the $x$
and $y$ components. This is also seen in the multiplicity of the eigenvalues,
which are now $\lambdaup_1=1-2\sp t\Delta^2_z$ (multiplicity one),
$\lambdaup_2=1-2\sp t\Delta^2_x$ (multiplicity two), and $\lambdaup_3=1$ (multiplicity three),
where $\Delta_x^2$ and $\Delta_z^2$ are given by equations~\eqref{eqn:Delta_x} and \eqref{eqn:Delta_z},
respectively. With respect to the Galactic plane, we decompose into transverse ($x$ and $y$)
and longitudinal components ($z$), labelled collectively by
`$\perp$' and `$\|$', respectively.

We now have two distinct $t$-dependent eigenvalues (where previously in the full-sky
case there was only one), thus producing one additional pole in the
moment generating function:
\be
    {M}_{Y}(t)={M}_{Y_\perp}(t)\ssp{M}_{Y_\|}(t),
\ee
with ${M}_{Y_\perp}(t)\equiv{M}_{Y_x}(t)\ssp{M}_{Y_y}(t)$,
${M}_{Y_x}(t)={M}_{Y_y}(t)=(1-2\sp t\Delta^2_x)^{-1/2}$,
and ${M}_{Y_\|}(t)={M}_{Y_z}(t)=(1-2\sp t\Delta^2_z)^{-1/2}$ (these can be found by considering the
eigenvalues of the covariance matrix~\eqref{eqn:cov-6x6-cutsky}).
The PDFs of each of these constituent moment generating functions are obtained by
the following inverse Laplace transforms:
%% \bea
%%     p_\perp(Y)
%%         &\equiv\frac{1}{2\pi\im}\int^{\gamma+\im\infty}_{\gamma-\im\infty} \frac{\rme^{tY}\dif t}{1+2\sp t\Delta^2_x}
%%         =\frac{\rme^{-Y/(2\Delta^2_x)}}{2\Delta^2_x}, \\[3.5pt]
%%     p_z(Y)
%%         &\equiv\frac{1}{2\pi\im}\int^{\gamma+\im\infty}_{\gamma-\im\infty} \frac{\rme^{tY}\dif t}{\big(1+2\sp t\Delta^2_z\big)^{1/2}}
%%         =\frac{\rme^{-Y/(2\Delta^2_z)}}{\big(2\pi\sp Y\Delta^2_z\big)^{1/2}}.
%% \eea
\begin{subequations}
\bea
    p_\perp(Y)
        &=\mathscr{L}^{-1}\big[{M}_{Y_\perp}(-t)\big](Y)=\frac{\rme^{-Y/(2\Delta^2_x)}}{2\sp\Delta^2_x}, \\
    p\Z{\|}(Y)
        &=\mathscr{L}^{-1}\big[{M}_{Y_\|}(-t)\big](Y)=\frac{\rme^{-Y/(2\Delta^2_z)}}{\big(2\sp\pi\sp Y\Delta^2_z\big)^{1/2}}.
\eea
\end{subequations}
By the convolution theorem of the unilateral Laplace transform, the PDF of the product
of moment generating functions $p(Y)=\mathscr{L}^{-1}[M_{Y_\perp}(-t)\ssp M_{Y_\|}(-t)]$ is
equal to the convolution $\mathscr{L}^{-1}[M_{Y_\perp}(-t)]\sp*\sp\mathscr{L}^{-1}[M_{Y_\|}(-t)]$,
which evaluates to
\bea
    p(Y)
        &=\int^Y_0\dif Y' p_\perp(Y-Y')\sp p\Z{\|}(Y') %\nonumber\\[-2pt]
        =\frac{1}{\sqrt{\sp\varepsilon}}\, {\rm erf}
            \bigg(\frac{\varepsilon Y}{2\Delta^2_z}\bigg)
            \frac{\rme^{-Y/(2\Delta^2_x)}}{2\Delta^2_x}, \nonumber
    % \label{eqn:p-Y-mask-lt0}
\eea
where $\varepsilon\equiv1-\Delta^2_z/\Delta^2_x$ gives the `eccentricity' of the ellipsoid spanned by the eigenvectors. The use of the error function in
the foregoing equation is valid for $\varepsilon>0$; for $\varepsilon<0$, which is the
case for the Galactic plane cut, we have instead
\bea
    p(Y)
        % &=\frac{1}{\sqrt{|\varepsilon|}}\ssp {\rm erfi}\bigg[\sqrt{|\varepsilon|}\sp\Big(\frac{Y}{2\Delta^2_z}\Big)^{1/2}\bigg]
        %     \frac{\rme^{-Y/(2\Delta^2_x)}}{2\Delta^2_x} \nonumber\\
        &=\frac{2}{\sqrt{\pi|\varepsilon|}}\ssp \rme^{|\varepsilon|Y/(2\Delta^2_z)}\sp
            F\bigg(\frac{\sqrt{|\varepsilon|}\sqrt{Y}}{\sqrt2\Delta_z}\bigg)\sp
            \frac{\rme^{-Y/(2\Delta^2_x)}}{2\Delta^2_x},
\eea
where $F(x)\equiv\rme^{-x^2}\int^x_0\dif t\,\rme^{t^2}$ is the Dawson function~\citep[][\href{https://dlmf.nist.gov/7.2.5}{7.2.5}]{DLMF}.

The probability distribution of the amplitude $\Dtot=\sqrt{\sp Y}$,
in the presence of a Galactic plane cut, is now obtained easily from ${p(\Dtot)=p(Y\!=\!\Dtot^2)\ssp|\sp\dif Y/\dif\Dtot|}$;
it is
%% \be
%%     p(X)
%%         =\frac{2X}{\sqrt{\pi|\varepsilon|}\sp\Delta^2_x}\ssp
%%             F\bigg(\sqrt{\frac{|\varepsilon|X^2}{2\Delta^2_z}}\bigg)\ssp
%%             \rme^{-X^2/(2\Delta^2_x)+|\varepsilon|X^2/(2\Delta^2_z)}.
%% \ee
\be\label{eqn:pD-cutsky-pre}
    p(\Dtot)\ssp\dif\Dtot
        =\frac{2\Dtot\ssp\dif\Dtot}{\sqrt{\pi|\varepsilon|}\sp\Delta^2_x}\ssp
            F\bigg(\frac{\sqrt{|\varepsilon|}\Dtot}{\sqrt2\Delta_z}\bigg)
            \exp\bigg(-\frac{\Dtot^2}{2\Delta^2_z}\bigg),
\ee
where we have used that $|\varepsilon|=\Delta^2_z/\Delta^2_x-1>0$, since for us $\varepsilon<0$.
We can bring equation~\eqref{eqn:pD-cutsky-pre} into a form
more closely resembling the full-sky case
(i.e.~the Maxwell--Boltzmann distribution)
by using $\Delta_x^2=\Delta_z^2/(1-\varepsilon)$ to rewrite it
in terms of $\Delta_z$ and $\varepsilon$.
After some algebra, and after defining the function
$\mathcal{F}(x)\equiv F(x)/x$, we arrive at our final form
[cf.~equation~\eqref{eqn:p-X}]:
\be\label{eqn:pD-cutsky}
\begin{split}
    p(\Dtot)\ssp\dif\Dtot
    &=(1-\varepsilon)\ssp     \mathcal{F}\bigg(\frac{\sqrt{|\varepsilon|}\Dtot}{\sqrt2\sp\Delta_z}\bigg)
     % &\quad\quad\times
     \frac{4}{\sqrt\pi}\sp\exp\bigg(\!-\frac{\Dtot^2}{2\Delta_z^2}\bigg)
     \frac{\Dtot\Zu2}{2\sp\Delta_z^2}\ssp\frac{\dif\Dtot}{\sqrt2\sp\Delta_z}.
\end{split}
\ee
This is equation~\eqref{eqn:p-Dtot-cutsky} in the main text.
As a consistency check, it is easy to verify using $\mathcal{F}(a\sqrt{\sp\varepsilon}\ssp)\to1$, as $\varepsilon\to0$, that the full-sky expression~\eqref{eqn:p-X}
derived in the previous section is recovered in the limit $\varepsilon\to0$.
This is the special case of statistical isotropy, for which $\Delta_x=\Delta_y=\Delta_z$.

The first three moments of the PDF~\eqref{eqn:pD-cutsky} are
\begin{subequations}
\bea
    &\langle\Dtot\rangle
        =\sqrt\frac{2}{\pi}
            \bigg[\Delta_z + \Delta_x\frac{{\rm arcsinh}\big(\sqrt{|\varepsilon|}\sp\big)}{\sqrt{|\varepsilon|}}\bigg], \\[2pt]
    &\langle\Dtot^2\rangle
        =\Delta^2_z + 2\Delta^2_x, \\[3pt]
        %% =(3-\varepsilon)\Delta^2_x, \\[3pt]
    &\langle\Dtot^3\rangle
        =\sqrt\frac{2}{\pi}\bigg[2\Delta^3_z+3\Delta_z\Delta^2_x
            +3\Delta^3_x\frac{{\rm arcsinh}\big(\sqrt{|\varepsilon|}\sp\big)}{\sqrt{|\varepsilon|}}\bigg].
\eea
\end{subequations}
In the limit $\varepsilon\to0$, these recover the usual moments of the Maxwell--Boltzmann distribution, i.e.\
$\langle\Dtot\rangle=\sqrt{\sp8/\pi}\sp\Delta$, $\langle\Dtot^2\rangle=3\Delta^2$, and
$\langle\Dtot^3\rangle=4\sqrt{\sp8/\pi}\sp\Delta^3$, with $\Delta=\Delta_x=\Delta_y=\Delta_z$.









\section{Dipole statistics on the cut sky}\label{app:galactic-cut}
In order to compare the dipole estimate with the theoretical value we will
convolve the full-sky predictions with the mask.

\subsection{Galactic plane cut\label{app:gal-cut}}
We make a {symmetric} cut about the Galactic plane, excluding a central region $|b|\leq b_{\rm cut}$.
Here $b_{\rm cut}=\ang{30}$, as in S21. This conservative cut removes half of the sky ($f_\mrm{sky}=0.5$). Since the
additional masks (accounting for poor-quality photometry, image artefacts, etc) covers merely $1.2\%$ of
the sky we will ignore their effect on the power, which principally affects the small angular
scales.
We will thus apply to the number-count field $\dif N/\dif\Omega(\n)$
the following mask:
\be
    \mathcal{W}(\vartheta)
        =
        \begin{cases}
            0, & \vartheta_-\leq\vartheta\leq\vartheta_+, \\
            1, & \text{otherwise},
        \end{cases}
\ee
where $\vartheta_\pm=\pi/2\pm b_{\rm cut}$. % and  $\vartheta_b=\pi/2+b_{\rm cut}=\pi-\vartheta_a$.
In terms of the Heaviside step function $\Theta$, this can also be written as
\be
    \mathcal{W}(\n)
        =\Theta(\n\cdot\hat{\mathbf{z}}-\mu\Z-)
        +\Theta(\mu\Z+-\n\cdot\hat{\mathbf{z}}),
\ee
with $\mu_{\pm}=\cos\vartheta_{\pm}$, and where
the first and second term selects the northern and southern caps, respectively.
Because $\mathcal{W}(\n)$ possesses azimuthal symmetry, its spherical harmonic expansion
is given purely in terms of the zonal modes ($m=0$):
$\mathcal{W}(\vartheta)=\sum_{\ell} \mathcal{W}_{\ell 0}\ssp Y_{\ell0}(\vartheta)$.
Furthermore, because $\mathcal{W}(\n)$ is parity symmetric,
$\mathcal{W}_{\ell m}=\int\dif^2\n~\mathcal{W}(\n)\sp Y^*_{\ell m}(\n)$ is equal to zero
for odd $\ell$;  for even $\ell$ these coefficients are non-zero and are
given by~\citep{Dahlen:2007sv}
\bea
    \mathcal{W}_{\ell m}
        % &\equiv\int\sin\vartheta\ssp\dif\vartheta\ssp\dif\phi\
        %     \mathcal{W}(\vartheta)\sp Y_{\ell m}^*(\vartheta,\phi) \nonumber\\
        &=2\sp\delta\Zu{\mrm{K}}_{m0}\ssp\sqrt{\frac{\pi}{2\ell+1}}\ssp
            \big[\mathcal{L}\Z{\ell-1}(\mu\Z-)-\mathcal{L}\Z{\ell+1}(\mu\Z-)\big]
         \quad\text{(even $\ell$)},
\eea
where $\calL_\ell(\mu)$ is the Legendre polynomial of degree $\ell$ [note $\mathcal{L}\Z{-1}(\mu)\equiv1$].
% (For the southern polar cap, $\mathcal{W}_{\ell m}\to -\mathcal{W}_{\ell m}$; for a double polar
% cap $\mathcal{W}_{\ell m}\to 2\mathcal{W}_{\ell m}$, and, because it is parity symmetric,
% $\mathcal{W}(\vartheta)=\mathcal{W}(\pi-\vartheta)$, only coefficients with even $\ell$ are non-zero.)
With $b_\mrm{cut}=\ang{30}$ we have for the first five non-zero coefficients
$\mathcal{W}_{00}\approx 1.77245$,
$\mathcal{W}_{20}\approx 1.48625$,
$\mathcal{W}_{40}\approx-0.623128$,
$\mathcal{W}_{60}\approx-0.131059$, and
$\mathcal{W}_{80}\approx 0.422182$. % Mathematica

\subsection{Mode coupling}\label{app:mode-coupling}
A cut on the sky breaks statistical
isotropy. In harmonic space this implies that the covariance
matrix of $a\Z{\ell m}$'s can have off-diagonal components, i.e.~we no longer have that
$\langle{a}\Z{\ell m}\sp {b}_{\ell'm'}^* \rangle=C_\ell\ssp\delta^\mrm{K}_{\ell\ell'}\sp\delta^\mrm{K}_{mm'}$.
Now, $a_{\ell m}$ with different $\ell$
or $m$ may be correlated.

We now apply the mask $\mathcal{W}(\n)$ to the number-count
field
$\dif N/\dif\Omega(\n)\to\mathcal{W}(\n)\dif N/\dif\Omega(\n)$.
Denoting by $\tilde{a}_{\ell m}$
and $\tilde{b}_{\ell m}$ the coefficients of the masked fields
$\mathcal{W}(\n)\sp\dqso(\n)$ and $\mathcal{W}(\n)(\dkin\cdot\n)$,
respectively, it can be shown using the properties of spherical
harmonics that we now have a linear coupling of the
underlying $a_{\ell m}$ and $b_{\ell m}$~\citep{1980lssu.book.....P}:
\be
    \tilde{a}\Z{\ell m}
        =\sum_{\ell' m'} \mathcal{A}\Z{\ell m\ell'm'}\,\sp a\Z{\ell'm'},\qquad
    \tilde{b}\Z{\ell m}
        =\sum_{\ell' m'} \mathcal{B}\Z{\ell m\ell'm'}\,\sp b\Z{\ell'm'}.
\ee
Here $\mathcal{A}_{\ell m\ell' m'}$ is a harmonic kernel which is
determined by the window function $\mathcal{W}(\n)$; for a general
$\mathcal{W}(\n)=\sum_{\ell,m} \mathcal{W}_{\ell m}Y_{\ell m}(\n)$, it
is given by~\citep{Hivon:2002jp}
\bea
    \mathcal{A}_{\ell_1 m_1 \ell_2 m_2}
        &\equiv\sum_{\ell_3,m_3}\mathcal{W}_{\ell_3 m_3}\int\dif^2\n\ Y_{\ell_1 m_1}(\n)\, Y\Zu{*}_{\ell_2m_2}(\n)\, Y_{\ell_3m_3}(\n) \nonumber\\
        &=(-1)^{m_2}\sum_{\ell_3,m_3}\sp \mathcal{W}\Z{\ell_3 m_3}
            \bigg[\frac{(2\ell_1+1)(2\ell_2+1)(2\ell_3+1)}{4\pi}\bigg]^{1/2}\! \nonumber\\[2pt]
        &\qquad\qquad\quad
            \times
            \begin{pmatrix}
                \ell_1 & \ell_2 & \ell_3 \\
                0 & 0 & 0
            \end{pmatrix}
            \begin{pmatrix}
                \ell_1 & \ell_2 & \ell_3 \\
                m_1 & -m_2 & m_3
            \end{pmatrix},
       \label{eqn:A-wigner}
\eea
where the arrays are Wigner $3j$ symbols,
% \footnote{\url{https://dlmf.nist.gov/34.3}; see also \citet*{Varshalovich_book}.}
coefficients that encode selection rules for the allowed combinations of $\ell$'s and $m$'s.
Using that $\mathcal{A}_{\ell m\ell'm'}=\mathcal{A}_{\ell m\ell'm}\ssp\delta^\mrm{K}_{mm'}$ (by azimuthal symmetry), and
$\mathcal{A}_{\ell m\ell'm'}^*=\mathcal{A}_{\ell m\ell'm'}$ (because $\mathcal{W}^*_{\ell0}=\mathcal{W}_{\ell0}$),
we have
% \bea
%     \big\langle\tilde{a}^{\phantom{*}}\Z{\ell_1m_1}\tilde{a}\Z{\ell_2m_2}^*\big\rangle
%         &=\sum_{\ell_1' \ell_2' m_1' m_2'} \mathcal{A}_{\ell_1m_1\ell_1'm_1'}\ssp \mathcal{A}_{\ell_2m_2\ell_2'm_2'}^*\ssp
%             \langle a_{\ell_1'm_1'}\sp a^*_{\ell_2'm_2'}\rangle \nonumber\\
%         &=\sum_{\ell_1' \ell_2' m_1' m_2'} \mathcal{A}_{\ell_1m_1\ell_1'm_1}\ssp \mathcal{A}_{\ell_2m_2\ell_2'm_2}\ssp
%             \delta_{m_1m_1'}\sp\delta_{m_2m_2'} \nonumber\\[-2pt]
%             &\qquad\qquad\quad \times C_{\ell_1'}\sp\delta_{\ell_1'\ell_2'}\sp\delta_{m_1'm_2'} \nonumber\\
%         &=\sum_{\ell_1' \ell_2' m_1'} \mathcal{A}_{\ell_1m_1\ell_1'm_1}\ssp \mathcal{A}_{\ell_2m_2\ell_2'm_2}\ssp
%             \delta_{m_1m_1'}\sp\delta_{m_2m_1'}\sp
%             C_{\ell_1'}\sp\delta_{\ell_1'\ell_2'} \nonumber\\
%         &=\sum_{\ell_1' m_1'} \mathcal{A}_{\ell_1m_1\ell_1'm_1}\ssp \mathcal{A}_{\ell_2m_2\ell_1'm_2}\ssp
%             \delta_{m_1m_1'}\sp\delta_{m_2m_1'}\sp
%             C_{\ell_1'} \nonumber\\
%         &=\sum_{\ell_1'} \mathcal{A}_{\ell_1m_1\ell_1'm_1}\ssp \mathcal{A}_{\ell_2m_2\ell_1'm_2}\ssp
%             \delta_{m_1m_2}\sp
%             C_{\ell_1'} \nonumber\\
%         &=\delta_{m_1m_2}\sum_{\ell} \mathcal{A}_{\ell_1m_1\ell m_1}\ssp \mathcal{A}_{\ell_2m_2\ell m_1}\sp C_{\ell}
% \eea
for two arbitrary fields, each with azimuthally-symmetric masks (which need not be identical),
\bea
    \big\langle\tilde{a}^{\phantom{*}}\Z{\ell_1m_1}\tilde{b}\Z{\ell_2m_2}^*\big\rangle
        &=\sum_{\ell_1',\ell_2',m_1',m_2'} \mathcal{A}_{\ell_1m_1\ell_1'm_1'}\ssp \mathcal{B}_{\ell_2m_2\ell_2'm_2'}^*\ssp
            \big\langle a_{\ell_1'm_1'}\sp b^*_{\ell_2'm_2'}\big\rangle \nonumber\\
        % &=\sum_{\ell_1' \ell_2' m_1' m_2'} \mathcal{A}_{\ell_1m_1\ell_1'm_1}\ssp \mathcal{B}_{\ell_2m_2\ell_2'm_2}\ssp
        %     \delta_{m_1m_1'}\sp\delta_{m_2m_2'} \nonumber\\[-2pt]
        %     &\qquad\qquad\quad \times C^{ab}_{\ell_1'}\sp\delta_{\ell_1'\ell_2'}\sp\delta_{m_1'm_2'} \nonumber\\
        % &=\sum_{\ell_1' \ell_2' m_1'} \mathcal{A}_{\ell_1m_1\ell_1'm_1}\ssp \mathcal{B}_{\ell_2m_2\ell_2'm_2}\ssp
        %     \delta_{m_1m_1'}\sp\delta_{m_2m_1'}\sp
        %     C^{ab}_{\ell_1'}\sp\delta_{\ell_1'\ell_2'} \nonumber\\
        % &=\sum_{\ell_1' m_1'} \mathcal{A}_{\ell_1m_1\ell_1'm_1}\ssp \mathcal{B}_{\ell_2m_2\ell_1'm_2}\ssp
        %     \delta_{m_1m_1'}\sp\delta_{m_2m_1'}\sp
        %     C^{ab}_{\ell_1'} \nonumber\\
        % &=\sum_{\ell_1'} \mathcal{A}_{\ell_1m_1\ell_1'm_1}\ssp \mathcal{B}_{\ell_2m_2\ell_1'm_2}\ssp
        %     \delta_{m_1m_2}\sp
        %     C^{ab}_{\ell_1'} \nonumber\\
        &=\delta^\mrm{K}_{m_1m_2}\sum_{\ell} \mathcal{A}_{\ell_1m_1\ell m_1}\ssp
            \mathcal{B}_{\ell_2m_2\ell m_1}\sp C_{\ell}.
\eea

In our case -- with $\tilde{a}_{\ell m}$ corresponding to the clustering field and $\tilde{b}_{\ell m}$
corresponding to the kinematic field -- we only need to consider $\ell_1=\ell_2=1$, with identical mask 
applied to both fields ($\mathcal{A}=\mathcal{B}$).
In terms of the underlying power spectra we have
\begin{subequations}\label{eqn:power-masked}
\bea
    \big\langle\tilde{a}^{\phantom{*}}\Z{1m}\tilde{b}\Z{1m'}^*\big\rangle
        &=\delta^\mrm{K}_{mm'}\sp|\mathcal{A}_{1m1m}|^2\sp C\Zu{\ssp\kin\clu}_1, \\[5pt]
    \big\langle\tilde{a}^{\phantom{*}}\Z{1m}\tilde{a}\Z{1m'}^*\big\rangle
        &=\delta^\mrm{K}_{mm'}\sum_{\ell\ \text{odd}} |\mathcal{A}_{1m\ell m}|^2 C\Zu{\ssp\clu\clu}_{\ell}, \\
    \big\langle\tilde{b}^{\phantom{*}}\Z{1m}\tilde{b}\Z{1m'}^*\big\rangle
        &=\delta^\mrm{K}_{mm'}\sp|\mathcal{A}_{1m1m}|^2\sp C\Zu{\ssp\kin\kin}_1,
\eea
\end{subequations}
where only odd $\ell$ need to be considered because of the selection rules of Wigner $3j$ symbols in equation~\eqref{eqn:A-wigner}.
Note that no off-diagonal entries are induced in the
covariance matrix in harmonic space. However, in Cartesian space the covariance matrix
is no longer proportional to the identity matrix (although it remains diagonal).
Though we will no longer have spherical symmetry we still have azimuthal symmetry.


For $b_\mrm{cut}=\ang{30}$, the first three coupling coefficients are
$|\mathcal{A}_{1010}|^2\approx 0.76563$,
$|\mathcal{A}_{1030}|^2\approx 0.046143$, and
$|\mathcal{A}_{1050}|^2\approx 0.034272$.
% $|\mathcal{A}_{1070}|^2\approx 0.0050754$, and
% $|\mathcal{A}_{1090}|^2\approx 0.0016420$.
Also
$|\mathcal{A}_{1111}|^2\approx 0.0976563$,
$|\mathcal{A}_{1131}|^2\approx 0.155731$, and
$|\mathcal{A}_{1151}|^2\approx 0.00955939$.
% $|\mathcal{A}_{1171}|^2\approx 0.00857833$, and
% $|\mathcal{A}_{1191}|^2\approx 0.0141864$.
These values indicate a substantial loss of power for the $x$- and $y$-components of the dipole.



% \begin{table}
% \centering
% \caption{Harmonic mode-coupling coefficients $\mathcal{A}_{1m\ell m}$ for galactic plane cuts of
% two different sizes. Note that for even $\ell$, $\mathcal{A}_{1m\ell m}=0$.}
% \setlength{\tabcolsep}{8pt}
% \begin{tabular}{ cllll }
% %% \begin{tabular}{ |c|l|l|l|l| }
% %% \hline
% %% \hline
% \toprule
% & \multicolumn{2}{c}{$b_\mrm{cut}=\ang{25}$} & \multicolumn{2}{c}{$b_\mrm{cut}=\ang{30}$} \\
% \cmidrule(lr){2-3}
% \cmidrule(lr){4-5}
% $\ell$ & \quad$m=0$ & \quad$m=\pm1$ & \quad$m=0$ & \quad$m=\pm1$ \\
% \midrule
% 1 & $\phantom{-}0.92452$ & $\phantom{-}0.40381$ & $\phantom{-}0.875$ & $\phantom{-}0.3125$ \\
% 3 & $\phantom{-}0.14206$ & $\phantom{-}0.40008$  & $\phantom{-}0.21481$ & $\phantom{-}0.39463$ \\
% 5 & $-0.15104$ & $-0.18404$ & $-0.18513$ & $-0.097772$ \\
% 7 & $\phantom{-}0.10447$ & $-0.0022393$ & $\phantom{-}0.071242$ & $-0.092619$ \\
% 9 & $-0.030757$ & $\phantom{-}0.10174$ & $\phantom{-}0.040522$ & $\phantom{-}0.11911$ \\
% 11& $-0.032269$ & $-0.10274$ & $-0.080434$ & $-0.034391$ \\
% 13& $\phantom{-}0.057766$ & $\phantom{-}0.039514$ & $\phantom{-}0.042809$ & $-0.054173$ \\
% %% \hline
% %% \hline
% \bottomrule
% \end{tabular}
% % and that because $\mathcal{A}_{1-1\ell-1}=\mathcal{A}_{11\ell1}$ we do not report the $m=-1$ values.}
% % Here we report five sig figs
% \end{table}





% The particular kernel for our mask is
% \bea
%     \mathcal{A}_{1 m \ell m}
%         &=(-1)^{m}\sum_{\ell'=\ell\pm1}\sp \mathcal{W}\Z{\ell'0}
%             \bigg[\frac{3(2\ell+1)(2\ell'+1)}{4\pi}\bigg]^{1/2} \nonumber\\[-1pt]
%             &\qquad\qquad\qquad
%             \times
%             \begin{pmatrix}
%                 1 & \ell & \ell' \\
%                 0 & 0 & 0
%             \end{pmatrix}
%             \begin{pmatrix}
%                 1 & \ell & \ell' \\
%                 m & -m & 0
%             \end{pmatrix}.
% \eea
% For a double polar cap window function we have $W_{\ell m}=0$ if $m\neq0$, and 
% \cite{Dahlen:2007sv}
% \bea
%     \mathcal{W}_{\ell 0}
%         &=\int\dif^2\n\ \mathcal{W}(\n)\sp Y_{\ell m}^*(\n) \nonumber\\
%         &=2
%         \begin{cases}
%             \sqrt{\dfrac{\pi}{2\ell+1}}\,
%             \big[\mathcal{L}\Z{\ell-1}(\mu_{\rm cut})
%                 -\mathcal{L}\Z{\ell+1}(\mu_{\rm cut})\big], & \text{$\ell$ is even}, \\[7pt]
%             0, & \text{$\ell$ is odd},
%         \end{cases}
% \eea
% with $\mathcal{L}\Z{-1}(\mu)\equiv1$ and
% $\mu_{\rm cut}=\cos(\pi/2-|b_{\rm cut}|)$. (The factor of two is because $\mathcal{W}(\n)$
% is boxcar function, i.e.~it consists of two Heaviside step functions; the odd moments
% are zero because of azimuthal symmetry in $\mathcal{W}(\n)$.)


\subsection{Dipole dispersions}\label{app:stats-cut}
We now need to relate the statistics in harmonic space back to the statistics
in Cartesian space. To facilitate this it will be
convenient to define the complex-valued vectors
$\bm{\tilde{a}}\Z1\equiv(\tilde{a}_{1-1},\tilde{a}_{11},\tilde{a}_{10})\Zu\T$
and $\bm{\tilde{b}}\Z1\equiv(\tilde{b}_{1-1},\tilde{b}_{11},\tilde{b}_{10})\Zu\T$.
We then have the following linear relation between the Cartesian- and harmonic-space dipoles:
\be
    \dint=\sqrt{\frac{3}{4\pi}}\,\mat{U}^\dagger\ssp\bm{\tilde{a}}\Z1, \quad
    \dkin=\sqrt{\frac{3}{4\pi}}\,\mat{U}^\dagger\ssp\bm{\tilde{b}}\Z1,
    %% \bm{\tilde{a}}\Z1=\sqrt{\frac{4\pi}{3}}\, \mat{U}\ssp\dint\,, \quad
    %% \bm{b}\Z1=\sqrt{\frac{4\pi}{3}}\, \mat{U}\ssp\dkin\,, \quad
    %% \mathbf{U}
    %% =
    %% \begin{pmatrix}
    %%     \!\phantom{-}1/\sqrt{\ssp2}    &   \im/\sqrt{\ssp2}  &   0   \\[2pt]
    %%     \!          -1/\sqrt{\ssp2}    &   \im/\sqrt{\ssp2}  &   0   \\[2pt]
    %%     \!0           &       0        &   1
    %% \end{pmatrix},
\ee
where the tilde indicates coefficients of masked fields, and $\mat{U}$ is given by equation~\eqref{eq:U}.
With these relations we can express the covariances of these vectors,
$\langle\mbf{d}_\mrm{X}^{\phantom\T}\mbf{d}_\mrm{Y}^\T\rangle$, in terms of the
underlying power spectra:
%% How does the analysis leading to equation~\eqref{eqn:p-Dtot-nocut} for $p(\Dtot)$ change with the
%% inclusion of this mask? This is best understood in terms of the covariance matrix, which
%% we recall can be written in block-matrix form with the $3\times3$ matrices
\begin{subequations}
\bea
    \big\langle\dkin\Zu{\phantom\T}\dkin\Zu\T\big\rangle
        %% &=\frac{3}{4\pi}\,\big\langle(\mathbf{U}^\dagger\bm{b}_1)(\mathbf{U}^\dagger\bm{b}_1)\Zu\dagger\big\rangle 
        &=\frac{3}{4\pi}\,\mathbf{U}^\dagger\langle\bm{\tilde{b}}_1^{\phantom\dagger}\bm{\tilde{b}}_1^\dagger\rangle\mathbf{U} \nonumber\\[-6pt]
        %% =\frac{3}{4\pi}\,{\rm diag}\big(\sp C^{bb}_1,\sp\,C^{bb}_1,\sp\,C^{bb}_1\big), \\[4pt]
        %% =\frac{3}{4\pi}\,C\Zu{\ssp\kin\kin}_1\,\mat{I}_3, \\[4pt]
        &={\rm diag}\big(\sp|\mathcal{A}_{1111}|^2,\sp\,|\mathcal{A}_{1111}|^2,\sp\,|\mathcal{A}_{1010}|^2\big)\sp\frac{3C_1\Zu{\ssp\kin\kin}}{4\pi}, \\
    \big\langle\dkin\Zu{\phantom\T}\dint\Zu\T\big\rangle
        %% &=\frac{3}{4\pi}\,\big\langle(\mathbf{U}^\dagger\bm{b}_1)(\mathbf{U}^\dagger\bm{\tilde{a}}_1)\Zu\dagger\big\rangle
        &=\frac{3}{4\pi}\,\mathbf{U}^\dagger\langle\bm{\tilde{b}}_1^{\phantom\dagger}\bm{\tilde{a}}_1^\dagger\rangle\mathbf{U} \nonumber\\[-6pt]
        &={\rm diag}\big(\sp|\mathcal{A}_{1111}|^2,\sp\,|\mathcal{A}_{1111}|^2,\sp\,|\mathcal{A}_{1010}|^2\big)\sp\frac{3C_1\Zu{\ssp\kin\clu}}{4\pi}, \\
    \big\langle\dint\Zu{\phantom\T}\dint\Zu\T\big\rangle
        %% &=\frac{3}{4\pi}\,\big\langle(\mathbf{U}^\dagger\bm{\tilde{a}}_1)(\mathbf{U}^\dagger\bm{\tilde{a}}_1)\Zu\dagger\big\rangle
        &=\frac{3}{4\pi}\,\mat{U}^\dagger\langle\bm{\tilde{a}}_1^{\phantom\dagger}\bm{\tilde{a}}_1^\dagger\rangle\mat{U} \nonumber\\[-4pt]
        &=\sum_{\ell\ \text{odd}}
         {\rm diag}\big(\sp|\mathcal{A}_{11\ell1}|^2,\sp\,|\mathcal{A}_{11\ell1}|^2,\sp\,|\mathcal{A}_{10\ell0}|^2\big)\sp\frac{3C_\ell\Zu{\ssp\clu\clu}}{4\pi},
\eea
\end{subequations}
The block-matrix form of the $6\times6$ cut-sky covariance~\eqref{eqn:cov-6x6-cutsky} is constructed
out of these $3\times3$ covariances.
The first matrix $\langle\dkin\Zu{\phantom\T}\dkin\Zu\T\rangle\propto\mathbf{I}_3$ is unchanged from
before; the second and third matrices are diagonal but have different dispersions from before. For instance,
the dispersion of the $z$-component is proportional to
$\langle|\tilde{a}\Z{10}|^2\rangle=\sum_\ell|\mathcal{A}_{10\ell0}|^2\sp C^{}_\ell$, while the $x$- and
$y$-components both have dispersions proportional to $\langle|\tilde{a}\Z{11}|^2\rangle=\sum_\ell|\mathcal{A}_{11\ell1}|^2\sp C^{}_\ell$,
with $\mathcal{A}_{11\ell1}\neq \mathcal{A}_{10\ell0}$.

We can relate the harmonic form to the dispersions in the cut-sky covariance matrix~\eqref{eqn:cov-6x6-cutsky}
as follows. Again, because of azimuthal symmetry, we recall that the dispersions for the $x$- and
$y$-components are equal (though modified from before),
e.g.~$\sigma^2_{\clu_x\clu_x}=\sigma^2_{\clu_y\clu_y}$, $\sigma^2_{\kin_x\clu_x}=\sigma^2_{\kin_y\clu_y}$, etc.
The new dispersions, in terms of the spectra, now read
% \begin{subequations}
% \bea
%     \sigma\Zu2\Z{\kin_x\kin_x}
%         &=\frac{3}{4\pi}\ssp|\mathcal{A}_{1111}|^2\ssp C\Zu{\ssp\kin\kin}_1,\quad\:
%     \sigma\Zu2\Z{\clu_x\clu_x}
%         =\frac{3}{4\pi}\sum_{\ell\ \text{odd}}|\mathcal{A}_{11\ell1}|^2\ssp C\Zu{\ssp\clu\clu}_\ell,\quad\:
%     \sigma\Zu2\Z{\kin_x\clu_x}
%         =\frac{3}{4\pi}\ssp|\mathcal{A}_{1111}|^2\ssp C\Zu{\ssp\kin\clu}_1, \\
%         %% =\mathcal{A}_{1111}\,\sigkinint^2, \\
%     \sigma\Zu2\Z{\kin_z\kin_z}
%         &=\frac{3}{4\pi}\ssp|\mathcal{A}_{1010}|^2\ssp C\Zu{\ssp\kin\kin}_1,\quad\:
%     \sigma\Zu2\Z{\clu_z\clu_z}
%         =\frac{3}{4\pi}\sum_{\ell\ \text{odd}}|\mathcal{A}_{10\ell0}|^2\ssp C\Zu{\ssp\clu\clu}_\ell,\quad\:
%     \sigma\Zu2\Z{\kin_z\clu_z}
%         =\frac{3}{4\pi}\ssp|\mathcal{A}_{1010}|^2\ssp C\Zu{\ssp\kin\clu}_1,
%         %% =\mathcal{A}_{1010}\,\sigkinint^2, 
% \eea
% \end{subequations}
\bea
    \sigma\Zu2\Z{\kin_x\kin_x}
        &=\frac{3}{4\pi}\ssp|\mathcal{A}_{1111}|^2\ssp C\Zu{\ssp\kin\kin}_1,
    &&\sigma\Zu2\Z{\kin_z\kin_z}
        =\frac{3}{4\pi}\ssp|\mathcal{A}_{1010}|^2\ssp C\Zu{\ssp\kin\kin}_1, \nonumber\\
    \sigma\Zu2\Z{\clu_x\clu_x}
        &=\frac{3}{4\pi}\sum_{\ell\ \text{odd}}|\mathcal{A}_{11\ell1}|^2\ssp C\Zu{\ssp\clu\clu}_\ell,
    &&\sigma\Zu2\Z{\clu_z\clu_z}
        =\frac{3}{4\pi}\sum_{\ell\ \text{odd}}|\mathcal{A}_{10\ell0}|^2\ssp C\Zu{\ssp\clu\clu}_\ell, \nonumber\\
    \sigma\Zu2\Z{\kin_x\clu_x}
        &=\frac{3}{4\pi}\ssp|\mathcal{A}_{1111}|^2\ssp C\Zu{\ssp\kin\clu}_1,
    &&\sigma\Zu2\Z{\kin_z\clu_z}
        =\frac{3}{4\pi}\ssp|\mathcal{A}_{1010}|^2\ssp C\Zu{\ssp\kin\clu}_1, \nonumber
\eea
where $C\Zu{\mrm{XY}}_\ell$ is given by equation~\eqref{eqn:Cab-ell}.
The dispersions related to the total dipole are given in terms of these by
equations~\eqref{eqn:Delta_x} and \eqref{eqn:Delta_z}.
For \LCDM, though we have $C\Zu{\ssp\clu\clu}_\ell>C\Zu{\ssp\clu\clu}_1$, for $\ell\geq2$, the
leakage of power into the dipole is not enough to overcome the loss of power
across all $\ell$, e.g.~$\sigma^2_{\kin_z\kin_z}=|A_{1010}|^2\sigkin^2\approx0.77\sigkin^2$, a
$23\%$ loss of kinematic dipole power.
Because of this we find that we only need to sum up the first few odd multipoles to attain
convergence (here we take $\ell_\mrm{max}=9$).




